\subsection{The abelian restrictions on the actual pure torsion frames}

\paragraph{D23.H10 (proved; independent whole-proof check).}
Let $h=(1,f)$ be an actual multiplier-one gentle element,
with the positive marking
\[
 B=C_2\amalg C_3=\langle s,r\rangle,\qquad
 F=\langle x,y\rangle,\quad x=(sr)^{-2},\quad y=sxs.
\]
Let its actual pure frames be
\[
       sf=\beta s\beta^{-1},\qquad
       f^{-1}r=A r A^{-1},\qquad A,\beta\in\widehat F.
\]
Then
\[
             A\in\widehat F',\qquad
       \beta^{\mathrm{ab}}=(-\rho,-\rho),\qquad
                      \rho\in2\widehat{\mathbb Z}.
\]
In particular, in any finite marked quotient whose pure
image has abelianization $C_{12}^2$ on $x,y$, a normalized
actual element with $\rho=0\pmod3$ must admit a conjugator
from the \emph{derived pure image} between $r$ and
$f^{-1}r$. Its harmonic conjugator between $s$ and $sf$
must have abelian coordinate either $(0,0)$ or $(6,6)$.
An arbitrary pure conjugator is insufficient for either
conclusion. These are separate conjugator requirements;
the same conjugator is not required to serve both torsion
generators.

\begin{proof}
The existence and uniqueness of the actual pure frames,
and the diagonal abelian form of $\beta$, are the supplied
LM7\_PUSH E.65.1. Rearranging the cubic frame equation gives
\[
                    f=(rAr^{-1})A^{-1}.
\]
The action $R$ of conjugation by $r$ on
$F^{\mathrm{ab}}=\widehat{\mathbb Z}^2$ is the standard
order-three peripheral rotation. Its characteristic
polynomial is $t^2+t+1$, so $\det(R-I)=3$.
Since $f\in\widehat F'$, abelianizing the last display gives
\[
                       (R-I)A^{\mathrm{ab}}=0.
\]
Multiply by the integral adjugate. Multiplication by $3$
is injective on $\widehat{\mathbb Z}^2=\prod_p\mathbb Z_p^2$,
so $A^{\mathrm{ab}}=0$. This proves $A\in\widehat F'$.
One must perform this argument in the torsion-free
profinite integral module: its reduction modulo $12$ alone
has nonzero order-three fixed vectors.

For clarity, the parity restriction on $\rho$ can be read
directly from E.65.11. In the free class-two quotient,
write
\[
          \beta^{-1}=x^\rho y^\rho[x,y]^d.
\]
The multiplier-one three-cycle kills the central
commutator coefficient of $f$, since three times that
coefficient is zero in $\widehat{\mathbb Z}$. The
interchange $\theta(x)=y$, $\theta(y)=x$ and
$f=\theta(\beta^{-1})^{-1}\beta^{-1}$ then give
\[
                            \rho^2+2d=0.
\]
Reduction in $\mathbb Z_2$ implies that $\rho$ is even;
at odd primes multiplication by $2$ is a unit, giving the
stated inclusion $\rho\in2\widehat{\mathbb Z}$.

A continuous surjection to a finite group sends the
closed commutator subgroup onto the finite derived
subgroup. Thus the cubic conjugator in the finite pure
image lies in its derived subgroup. Its harmonic
abelian coordinates are the reduction of
$(-\rho,-\rho)$. If additionally $\rho=0\pmod3$,
the only possibilities modulo $12$ are $\rho=0,6$,
giving precisely the two displayed harmonic cosets.
\end{proof}

\paragraph{Source and scope.}
The exact supplied inputs are LM7\_PUSH E.65.1 and the
paragraph \emph{The integral sign that fixes the lift},
especially E.65.11, in
\path{sources/LM7_PUSH/E_cm_character_agent.tex}.
This argument does not identify the full character
$\rho$ with a Kummer character. Supplied A.102 gives only
the congruence modulo three needed for D23.H2; no stronger
arithmetic normalization is inserted here.
